\subsection{JSON Data Boundaries}

JSON is a text-based data interchange format. It is not Python syntax, even though many JSON structures look similar to Python dictionaries and lists. JSON exists as text outside the Python runtime. When Python reads JSON, it parses that text and creates Python objects. When Python writes JSON, it serializes Python objects back into JSON text.

This makes JSON another boundary-crossing mechanism. The program does not merely read characters; it reconstructs a structured tree of values.

\subsubsection{JSON as Text-Based Tree Serialization: Objects, Arrays, Strings, Numbers, Booleans, and Null}

A JSON document represents data as a tree. The main JSON value categories are objects, arrays, strings, numbers, booleans, and null.

A JSON object is written with braces and contains name-value pairs. A JSON array is written with square brackets and contains ordered values. JSON strings use double quotes. JSON booleans are written as \texttt{true} and \texttt{false}. JSON null is written as \texttt{null}.

\begin{verbatim}
json_text = """
{
  "name": "Chuck Norris",
  "level": 42,
  "active": true,
  "quote": "Chuck Norris can parse invalid JSON.",
  "tags": ["joke", "fact", "runtime"],
  "metadata": {
    "source": "course demo",
    "verified": false
  },
  "missing_value": null
}
"""

print("JSON text:")
print(json_text)

print("Runtime object before parsing:")
print(f"type(json_text) = {type(json_text)}")
print(f"len(json_text)  = {len(json_text)}")

print()
print("Selected names from dir(json_text):")
for name in ["strip", "splitlines", "replace", "startswith", "encode"]:
    print(f"{name:12} -> {name in dir(json_text)}")
\end{verbatim}

At this point, \texttt{json\_text} is only a Python string. Python has not yet treated it as structured data. The braces, brackets, colons, commas, and quoted names are characters inside a \texttt{str} object.

The JSON parser is what converts that text into Python dictionaries, lists, strings, integers, booleans, and \texttt{None}.

\begin{verbatim}
import json

json_text = """
{
  "name": "Chuck Norris",
  "level": 42,
  "active": true,
  "quote": "Chuck Norris can parse invalid JSON.",
  "tags": ["joke", "fact", "runtime"],
  "metadata": {
    "source": "course demo",
    "verified": false
  },
  "missing_value": null
}
"""

data = json.loads(json_text)

print("Parsed Python object:")
print(data)

print()
print("Object types after parsing:")
print(f"type(data)                  = {type(data)}")
print(f"type(data['name'])          = {type(data['name'])}")
print(f"type(data['level'])         = {type(data['level'])}")
print(f"type(data['active'])        = {type(data['active'])}")
print(f"type(data['tags'])          = {type(data['tags'])}")
print(f"type(data['metadata'])      = {type(data['metadata'])}")
print(f"type(data['missing_value']) = {type(data['missing_value'])}")

print()
print("Selected names from dir(data):")
for name in ["keys", "values", "items", "get", "update", "pop"]:
    print(f"{name:10} -> {name in dir(data)}")

print()
print("Dictionary access:")
print(f"name          = {data['name']}")
print(f"level         = {data['level']}")
print(f"first tag     = {data['tags'][0]}")
print(f"metadata src  = {data['metadata']['source']}")
\end{verbatim}

The external JSON text has now become normal Python runtime data. The root object is a \texttt{dict}. The \texttt{tags} value is a \texttt{list}. The nested \texttt{metadata} value is another \texttt{dict}. JSON \texttt{true} became Python \texttt{True}. JSON \texttt{false} became Python \texttt{False}. JSON \texttt{null} became Python \texttt{None}.

This is serialization and deserialization. Serialization turns runtime structure into a stored or transferable textual representation. Deserialization parses that external representation and rebuilds runtime structure.

\subsubsection{Loading JSON into Python Dictionaries and Lists with \texttt{json.load()} and \texttt{json.loads()}}

The \texttt{json.loads()} function parses JSON from an existing Python string. The final \texttt{s} means string. The input is already loaded into memory as text.

\begin{verbatim}
import json

json_text = """
{
  "show": "The Simpsons",
  "line": "D'oh!",
  "episode_count": 42,
  "characters": ["Homer", "Marge", "Bart"],
  "still_running": true
}
"""

data = json.loads(json_text)

print("Loaded with json.loads():")
print(data)

print()
print("Object inspection:")
print(f"type(json_text) = {type(json_text)}")
print(f"type(data)      = {type(data)}")

print()
print("Selected names from dir(data):")
for name in ["keys", "items", "get", "setdefault", "pop"]:
    print(f"{name:12} -> {name in dir(data)}")

print()
print("Accessing parsed values:")
print(f"show          = {data['show']}")
print(f"line          = {data['line']}")
print(f"episode_count = {data['episode_count']}")
print(f"first character = {data['characters'][0]}")
\end{verbatim}

The \texttt{json.load()} function parses JSON from an open file object. It reads from the file object and returns the corresponding Python structure.

\begin{verbatim}
import json

path = "json_load_demo.json"

json_text = """
{
  "title": "Monty Python record",
  "answer": 42,
  "quotes": [
    "It's just a flesh wound.",
    "Nobody expects the Spanish Inquisition!"
  ],
  "valid": true
}
"""

with open(path, "w", encoding="utf-8") as file_obj:
    file_obj.write(json_text)

with open(path, "r", encoding="utf-8") as file_obj:
    data = json.load(file_obj)

    print("File object used by json.load():")
    print(f"file_obj       = {file_obj}")
    print(f"type(file_obj) = {type(file_obj)}")

print()
print("Loaded Python object:")
print(data)

print()
print("Object types:")
print(f"type(data)           = {type(data)}")
print(f"type(data['quotes']) = {type(data['quotes'])}")
print(f"type(data['answer']) = {type(data['answer'])}")
print(f"type(data['valid'])  = {type(data['valid'])}")

print()
print("Iterating through the quote list:")
for index, quote in enumerate(data["quotes"], start=1):
    print(f"{index:02d}: {quote}")
    print(f"    type(quote) = {type(quote)}")
\end{verbatim}

The difference is only the source of the JSON text.

\begin{itemize}
    \item \texttt{json.loads(json\_text)} parses JSON from a \texttt{str}.
    \item \texttt{json.load(file\_obj)} parses JSON from a text file object.
\end{itemize}

Both return Python objects. The root object may be a dictionary, a list, a string, a number, a Boolean, or \texttt{None}, depending on the JSON document.

A JSON document does not have to start with an object. It may also start with an array.

\begin{verbatim}
import json

json_text = """
[
  {"name": "Jerry", "score": 38},
  {"name": "Elaine", "score": 42},
  {"name": "George", "score": 17}
]
"""

rows = json.loads(json_text)

print("Root JSON value became a Python list:")
print(rows)

print()
print("Object inspection:")
print(f"type(rows) = {type(rows)}")
print(f"len(rows)  = {len(rows)}")

print()
print("Selected names from dir(rows):")
for name in ["append", "extend", "insert", "pop", "sort", "reverse"]:
    print(f"{name:10} -> {name in dir(rows)}")

print()
print("Inspecting list elements:")
for row in rows:
    print(f"name = {row['name']:8} | score = {row['score']:02d}")
    print(f"    type(row) = {type(row)}")
\end{verbatim}

Here the root Python object is a \texttt{list}, and each element of that list is a \texttt{dict}. JSON therefore represents nested trees, not only flat records.

\subsubsection{Writing Python Structures Back to JSON with \texttt{json.dump()} and \texttt{json.dumps()}}

The reverse operation converts Python structures into JSON text. The \texttt{json.dumps()} function serializes a Python object into a JSON string. The final \texttt{s} again means string.

\begin{verbatim}
import json

data = {
    "course": "Python I/O",
    "topic": "JSON",
    "answer": 42,
    "quotes": [
        "Serenity now!",
        "No soup for you!"
    ],
    "active": True,
    "missing": None
}

json_text = json.dumps(data, indent=2)

print("Original Python object:")
print(data)

print()
print("Serialized JSON text:")
print(json_text)

print()
print("Object types:")
print(f"type(data)      = {type(data)}")
print(f"type(json_text) = {type(json_text)}")

print()
print("Selected names from dir(json_text):")
for name in ["splitlines", "encode", "replace", "startswith", "count"]:
    print(f"{name:12} -> {name in dir(json_text)}")
\end{verbatim}

The variable \texttt{data} refers to a Python dictionary. The variable \texttt{json\_text} refers to a Python string containing JSON-formatted text. That string can be printed, stored in a file, sent over a network, or parsed later.

The \texttt{indent=2} argument asks the JSON serializer to format the output across multiple lines with indentation. Without it, the output is more compact.

\begin{verbatim}
import json

data = {
    "name": "Homer",
    "line": "D'oh!",
    "number": 42
}

compact_json = json.dumps(data)
pretty_json = json.dumps(data, indent=2)

print("Compact JSON:")
print(compact_json)

print()
print("Pretty JSON:")
print(pretty_json)

print()
print("Lengths:")
print(f"len(compact_json) = {len(compact_json)}")
print(f"len(pretty_json)  = {len(pretty_json)}")
\end{verbatim}

The \texttt{json.dump()} function serializes a Python object directly into an open text file object.

\begin{verbatim}
import json

path = "json_dump_demo.json"

data = {
    "name": "Arthur Dent",
    "answer": 42,
    "items": ["towel", "tea", "guide"],
    "panicking": False
}

with open(path, "w", encoding="utf-8") as file_obj:
    result = json.dump(data, file_obj, indent=2)

    print("File object used by json.dump():")
    print(f"file_obj       = {file_obj}")
    print(f"type(file_obj) = {type(file_obj)}")

print()
print("Return value from json.dump():")
print(f"result = {result}")

with open(path, "r", encoding="utf-8") as file_obj:
    content = file_obj.read()

print()
print("JSON file content read back as text:")
print(content)

print()
print("Type of file content:")
print(f"type(content) = {type(content)}")
\end{verbatim}

The function \texttt{json.dump()} writes into the file and returns \texttt{None}. The JSON file is still a text file. Reading it with \texttt{file\_obj.read()} returns a string. Reading it with \texttt{json.load()} returns parsed Python objects.

\begin{verbatim}
import json

path = "json_dump_demo.json"

with open(path, "r", encoding="utf-8") as file_obj:
    parsed_data = json.load(file_obj)

print("Parsed data:")
print(parsed_data)

print()
print("Object type:")
print(f"type(parsed_data) = {type(parsed_data)}")

print()
print("Accessing fields:")
print(f"name   = {parsed_data['name']}")
print(f"answer = {parsed_data['answer']}")
print(f"items  = {parsed_data['items']}")
\end{verbatim}

The practical distinction is:

\begin{itemize}
    \item \texttt{json.dumps(data)} returns a JSON \texttt{str};
    \item \texttt{json.dump(data, file\_obj)} writes JSON text to a file object;
    \item \texttt{json.loads(json\_text)} parses a JSON \texttt{str};
    \item \texttt{json.load(file\_obj)} parses JSON text from a file object.
\end{itemize}

\subsubsection{JSON Type Mapping Boundaries: \texttt{None} vs. \texttt{null}, \texttt{dict} vs. Object, \texttt{list} vs. Array}

JSON and Python do not use the same names for every value category. The parser maps JSON types into Python types, and the serializer maps Python types back into JSON types.

The most common mappings are:

\begin{itemize}
    \item JSON object maps to Python \texttt{dict};
    \item JSON array maps to Python \texttt{list};
    \item JSON string maps to Python \texttt{str};
    \item JSON number maps to Python \texttt{int} or \texttt{float};
    \item JSON \texttt{true} maps to Python \texttt{True};
    \item JSON \texttt{false} maps to Python \texttt{False};
    \item JSON \texttt{null} maps to Python \texttt{None}.
\end{itemize}

\begin{verbatim}
import json

json_text = """
{
  "dict_like": {"a": 1, "b": 2},
  "list_like": [10, 20, 42],
  "text_like": "D'oh!",
  "int_like": 42,
  "float_like": 3.14,
  "true_like": true,
  "false_like": false,
  "null_like": null
}
"""

data = json.loads(json_text)

print("Parsed values and Python types:")
for key, value in data.items():
    print(f"{key:12} -> {value!r:20} | type = {type(value)}")
\end{verbatim}

The mapping is visible in the printed types. There is no Python value named \texttt{null}. Python uses \texttt{None}. There are no Python Boolean literals named \texttt{true} and \texttt{false}. Python uses \texttt{True} and \texttt{False}.

The reverse mapping appears during serialization.

\begin{verbatim}
import json

data = {
    "python_none": None,
    "python_true": True,
    "python_false": False,
    "python_list": [1, 2, 42],
    "python_dict": {"line": "No soup for you!"}
}

json_text = json.dumps(data, indent=2)

print("Python object:")
print(data)

print()
print("JSON text:")
print(json_text)
\end{verbatim}

The Python value \texttt{None} is written as JSON \texttt{null}. Python \texttt{True} and \texttt{False} are written as JSON \texttt{true} and \texttt{false}. Python dictionaries become JSON objects, and Python lists become JSON arrays.

Not every Python object can be serialized to JSON automatically. JSON has a small type system. Arbitrary Python objects, file objects, sets, and many other runtime objects do not have a default JSON representation.

\begin{verbatim}
import json

data = {
    "normal_list": [1, 2, 3],
    "problem_set": {1, 2, 3}
}

print("Trying to serialize a set:")

try:
    json_text = json.dumps(data)
except TypeError as err:
    print("Serialization failed.")
    print(f"type(err) = {type(err)}")
    print(f"err       = {err}")

    print()
    print("Selected names from dir(err):")
    for name in ["args", "with_traceback", "add_note"]:
        print(f"{name:16} -> {name in dir(err)}")
\end{verbatim}

A \texttt{set} is a normal Python object, but JSON has no direct set type. The programmer must choose a representation, often a list.

\begin{verbatim}
import json

data = {
    "normal_list": [1, 2, 3],
    "set_as_list": list({1, 2, 3})
}

json_text = json.dumps(data, indent=2)

print("Serializable version:")
print(json_text)
\end{verbatim}

There is also a syntax boundary. JSON is not Python literal syntax. JSON strings require double quotes. JSON booleans are lowercase. JSON null is lowercase. Trailing commas are not allowed.

\begin{verbatim}
import json

bad_json = """
{
  "name": "Homer",
  "line": "D'oh!",
}
"""

print("Trying to parse invalid JSON:")

try:
    data = json.loads(bad_json)
except json.JSONDecodeError as err:
    print("JSON parsing failed.")
    print(f"type(err) = {type(err)}")
    print(f"err       = {err}")

    print()
    print("Selected names from dir(err):")
    for name in ["msg", "doc", "pos", "lineno", "colno"]:
        print(f"{name:8} -> {name in dir(err)}")

    print()
    print(f"line number   = {err.lineno}")
    print(f"column number = {err.colno}")
    print(f"message       = {err.msg}")
\end{verbatim}

The practical rule is strict: JSON is an external text format with its own grammar and type system. Python can map JSON values into Python objects and back again, but the boundary is real. A JSON document is not a dictionary until it is parsed, and a Python dictionary is not JSON until it is serialized.